% lit_transport.tex -- Literature 2.1: cooperative transport and link models (drafted 12 Oct 2026 from the
% thesis plan of 11 Oct, section 4 "2 Literature"). Condensed from modelling.tex and control.tex, which are no
% longer input (the distributed detail and the architecture figure are in Appendix F). Entries the plan lists for
% this section that are not yet in the bibliography (Lee/Sreenath/Kumar 2013, Mellinger 2010, Masone
% 2016, Geng 2020, Wang 2026) wait for the approval batch of plan section 7 and are marked in brackets.

\subsection{Cooperative transport and link models}
\label{sec:lit_models}

Several quadcopters carrying one suspended payload raise the mass a team can lift, constrain the payload's
swing through several attachment points and can control its full pose, at the price of a coupled system in
which every vehicle feels the others through the load \cite{tudelft, mohammadiDecentralizedMotionControl2018}.
The field is reviewed most recently by Lv et al.\ \cite{lv2026review} and, for suspended cables, by Estevez et
al.\ \cite{estevezReviewAerialTransportation2024}: most of it concerns trajectory tracking, tension allocation,
formation keeping and disturbance rejection for a team whose membership never changes, and most of it is
validated in simulation.

\paragraph{Link models.} The assumptions made about the link set both the size of the model and the regime in
which the controller is valid. The common choice is the rigid massless rod: a link of fixed length along a unit
vector, carrying a tension $\lambda_i \ge 0$ along it, so that the link can pull but not push
\cite{efficient, wahbaKinodynamicMotionPlanning2024}. It keeps the state compact and the equations smooth, and
it is exact for a taut cable under a hovering vehicle, which Michael et al.\ \cite{michael2011cooperative} treat
as a rod on ball joints; the non-negativity of the tension is then a modelling choice about when the link is
believed, not a property of the hardware \cite{six2018flying}. The mass-spring-damper chain of
De Carli et al.\ \cite{distributedNMPC} discretises the cable into lumped masses joined by spring-damper
elements, captures sag and oscillation, and costs six states per element per cable. Whatever the model, nearly
every controller assumes the links taut throughout; a slack link adds a unilateral constraint and a
discontinuity, and handling it rigorously is still open \cite{efficient}. The rig of this thesis carries its
ring on rigid rods, so the rod model is the hardware's own, except for the newcomer's tether
(Section~\ref{sec:overview}).

\paragraph{Centralised control.} Architectures are classified by their information structure, which sets the
computation, the scaling with team size and the robustness to a lost link \cite{nathanReviewAutonomousMultiagent2011}.
A central node with the global state can impose tension limits, collision avoidance and a force allocation
directly, at the cost of a solve that grows with the team and a single point of failure
\cite{dimakosStudyCentralisedDecentralised2021}. The benchmark for this thesis is Sun et al.\ \cite{tudelft}, who
solve a centralised optimal control problem over the full load--cable--vehicle dynamics at 10\,Hz, with thrust
limits, cable tautness and collision avoidance as path constraints, and track the planned trajectories onboard
with incremental nonlinear dynamic inversion (INDI), which compensates the cable force from inertial
measurements so that the load carries no sensor. Three and four vehicles reach 5\,m/s and 8\,m/s$^2$, at least
eight times the agility of earlier force-based methods. Two things in that work are left open and matter here:
the solve time grows exponentially with the team, and the results do not separate what the INDI force
compensation contributes from what the planner does. Full-pose nonlinear MPC scales to ten vehicles in
simulation at under 100\,ms \cite{sunNonlinearMPCFullPose2023}, or lowers its cost by parameterising the problem
on the payload's states alone \cite{liNonlinearModelPredictive2023}. Allocation methods separate the split of a
desired payload wrench into tensions: a cascade of quadratic programs after geometric control
\cite{efficient}, which scales about quadratically, or sequential quadratic programming on the squared tensions
\cite{vazquez-soquiSQPBasedCableTensionAllocation2026}. Kinodynamic planning finds a feasible path through
clutter offline and then enforces the full dynamics \cite{wahbaKinodynamicMotionPlanning2024}; leader-follower
control makes the payload the leader and splits the resultant tension among the vehicles, flown on hardware with
a point-mass payload \cite{tensionbased}; adaptive and barrier-function schemes guarantee constraints under
model uncertainty in simulation \cite{9945662}. All of these hold the team fixed.

\paragraph{Distributed and decentralised control.} Without a central node each vehicle solves for itself from
what it measures or what its neighbours send, which removes the single point of failure and the growth of one
solve with the team, and makes global objectives such as an exact force allocation harder
\cite{scattoliniArchitecturesDistributedHierarchical2009}. Distributed nonlinear MPC reaches centimetre-level
tracking with each vehicle solving its part of the global problem \cite{distributedNMPC}; event-triggered
variants update only when the system leaves its plan \cite{tasoojiEventTriggeredNonlinearModel,
tasoojiCooperativeControlMultiQuadrotors}; sliding-mode, adaptive, vision-based and learned controllers cover
the rest of the design space \cite{DFHSMC, arabPlanningDistributedControl2021, 9376103,
liMultiUAVCooperativeTransportation2021}. The line of work closest in spirit to a changing team is the passivity-based
control of Quan et al.\ \cite{quan2025swarm}: each vehicle damps its motion relative to the payload from local
measurements, with no communication, so that any disturbance, a member leaving included, injects energy that the
damping removes; its peer-reviewed form, a virtual-passivity network of nodes, springs and dampers, carries a
6\,kg load with six vehicles in field tests \cite{liu2026passivity}, and a virtual-tube extension steers the
formation through constrained spaces \cite{liuVirtualTubeBasedCooperativeTransport2026}. Appendix~\ref{sec:app_network}
reviews this family in more detail.

\paragraph{What this thesis takes.} The project plan also asked whether such a network or the trajectory-based
planner better balances robustness and tracking; that was answered on an earlier simulation plant and is recorded
as history in Appendix~\ref{sec:app_network}. What this thesis takes from the literature is the optimal control
problem of Sun et al.\ as its planner, with the team size and the attachment geometry made runtime parameters
(Section~\ref{sec:model}), and a tracking controller per vehicle on the ground in place of INDI, because the
flight controllers are unmodified (Section~\ref{sec:lit_execution}).
